\section{The Poisson score}
\label{sec:predictor-score}

We now construct the set promised by the predictor in the small-cell case.
The training set and its assignment cells come from Lemma~\ref{lem:training}.
The geometric estimates control dependence between tests through the same
query point.  Our immediate task is to turn those estimates into a score
which accepts many training points and few ambient points.

Throughout this section write
\begin{gather*}
 n=|S|=q^{1+v},\qquad n'=|S'|\in[n/4,n],\qquad B=q^d/n',\\
 \xi=L/100,\qquad \chi=P/100,\qquad \chi_0=P/10000.
\end{gather*}
We are in the case \(n>100qP\); the enumeration case was treated earlier.
For a query point \(x\), let \(O_x\subseteq S'\) be its own set: the union
of the assignment cells of the earliest flats containing \(x\) in the two
captured lists.  An absent list contributes the empty set.  Its fraction
\(f=|O_x|/n'\) is at most \(.08\).  The cell table determines this fraction,
including the intersection of the two cells.  It also determines whether a
sampled training point belongs to \(O_x\).

Let \(\mathscr C\) contain the cells from each list and their cross
intersections.  These are three partition systems, so
\(\sum_{C\in\mathscr C}|C|/n'\le3\).  For a hyperplane \(H\), put
\[
 \lambda_H=\frac{q|S'\cap H|}{n'},\qquad
 f_C=\frac{|C|}{n'},\qquad
 \lambda_{H,C}=\frac{q|C\cap H|}{n'}.
\]
The exceptional hyperplanes form the deterministic set
\begin{equation}
 \mathcal E=\{H:|\lambda_H-1|>.1\}
 \ \cup\!\bigcup_{C\in\mathscr C}
       \{H:|\lambda_{H,C}-f_C|>.02\}.
 \label{eq:exceptional-hyperplanes}
\end{equation}
A hyperplane is \emph{typical} if it does not belong to $\mathcal E$.
The encoder may use this set in analyzing a row; the decoder will not need
it to compute the score.

A \emph{pencil} at $x$ is the family of all hyperplanes containing $x$.
Take $R$ independent Poisson batches, each of mean $Lq$, whose points
are independent and uniform in $S'$.  Let $Z\subseteq\mathcal E$ be the
hyperplanes empty in every batch, put $z=|Z|$, and let $z_x$ count the
members of $Z$ containing $x$.

For each point \(x\), set \(b_0=e^{-L(1-f)}\).  Its score is
\begin{equation}
 A_x=\sum_{H\ni x}\prod_{r=1}^R
 \mathbf1\{\text{batch }r\text{ has no point of }O_x\cap H\}
 \bigl(\mathbf1\{\text{batch }r\text{ has no point of }(S'\setminus O_x)
                                      \cap H\}-b_0\bigr).
 \label{eq:poisson-score}
\end{equation}
The procedure includes every sampled point.  It includes an unsampled
point exactly when \(A_x<.5z/q\).  The samples and cell data determine
\eqref{eq:poisson-score}; only the integer \(z\) is additional information.

The score will approximate $z_x$: empty exceptional hyperplanes give
its main contribution, and centering suppresses typical hyperplanes.
We will show that $z_x$ is small for most training points and large
for most ambient points.

\begin{lemma}[Score procedure]
\label{lem:score-procedure}
Assume the small-cell conclusions of Lemma~\ref{lem:training}
and, when \(d\ge3\), the estimates of
Lemma~\ref{lem:overlap-geometry} for outside-own radial lines and hyperplane
pairs.  Assume also
\[
 |S||T|\ge q^{d+1}e^{-b},\qquad
 \frac{e(S,T)}{|S||T|}\le\frac{\tau}{q},\qquad
 b=O(D\sigma^{6\beta}),
\]
with the predictor's parameter ranges and \(2\le d\le4\).
Take \(R\) independent Poisson batches, each of mean \(Lq\), sampled
uniformly from \(S'\).  From the batches, the transmitted cell data and
one additional integer \(z\), this score procedure computes a set
\(W\) of points on the \(S\)-side.  With probability at least
\(\exp[-C(P\tau+1)]\), the row has at most \(2qP\) samples and satisfies
\[
 |W|\le n e^{CP},\qquad |W\cap S'|\ge c n'.
\]
These are simultaneous assertions under the original Poisson sampling
law.  The procedure does not condition on finding a successful row.
\end{lemma}

Here are the geometric inputs used in dimensions three and four, stated
explicitly to fix the interface.  Apart from \(o(n)\) centers, the number of radial lines
of outside-own strength at least \(a\) is at most
\(q^{2-2v}e^\chi a^{-100}\).  Among any pencil family whose outside-own
mass is at most two, the overlap graph at level \(a\) has ordered edge
count and maximum degree at most
\begin{equation}
 E_a\le B^2e^{2\chi}a^{-200},\qquad
 \Delta_a\le Be^{2\chi}a^{-200}.
 \label{eq:score-geometric-inputs}
\end{equation}
For the ambient analysis, after excluding an additional \(ne^{O(P)}\)
centers, its degree at \(t=1/(100p)\) is at most \(Be^{-3P}\), where \(p\)
is an even integer in \([10000\sigma/P,10000\sigma/P+2]\).
In the case \(g_1=(v-1)\sigma>\chi_0\), which occurs only in dimension
four, outside-own strengths are at most \(Ce^{-g_1}\).  The global number
of center--line incidences of strength at least \(a>e^{-.05g_0}\), where
\(g_0=(v-.5)\sigma\), is at most \(Cq^2/a^2\).
Lemma~\ref{lem:overlap-geometry} proves these statements without sampling.
All dyadic levels below range from \(q^{-10}\) to two.  Every positive
strength or overlap is at least \(q/n'\), so this range includes all
nonzero quantities which occur.

We shall repeatedly use the following consequences of the parameter
hierarchy:
\begin{equation}
 \log\sigma=o(P),\quad R\log L=o(P),\quad
 b+P\tau=o(L),\quad L=o(P),\quad P=o(\sigma).
 \label{eq:score-parameter-slack}
\end{equation}
They hold uniformly in the allowed range of \(D\), for fixed \(\eta>0\)
and sufficiently large \(q\).  In particular \(R/2>200\) and \(R>5000\).

\subsection{Empty exceptional hyperplanes}
\label{subsec:score-exceptional}

The incidence variance identity gives
\begin{equation}
 \sum_H(\lambda_{H,C}-f_C)^2\le CqBf_C,
 \qquad \sum_H(\lambda_H-1)^2\le CqB.
 \label{eq:cell-variance}
\end{equation}
Replacing the exact means by \(f_C\) and one changes only a negligible
term, since \(qp_d=1-1/Q_d\).  Summing the first estimate over the three
partition systems, then using the thresholds in
\eqref{eq:exceptional-hyperplanes}, gives
\begin{equation}
 |\mathcal E|\le CqB,\qquad
 \sum_{H\in\mathcal E}\lambda_H\le CqB.
 \label{eq:exceptional-mass}
\end{equation}
For the second assertion, use Cauchy--Schwarz and the full-set variance
in \eqref{eq:cell-variance}.

We call a pencil regular if it contains at most \(Be^\xi\) members of
\(\mathcal E\), and if the pencil sum of each clipped squared deviation
in \eqref{eq:cell-variance} is at most \(Be^\xi\). Here clipping replaces
a squared deviation $u^2$ by $\min\{1,u^2\}$.  Only \(o(n)\) centers
have an irregular pencil.  Indeed, a clipped weight \(w_H\in[0,1]\) has
\(\sum_Hw_H\le CqBf_C\) and \(\sum_Hw_H^2\le\sum_Hw_H\).  Applying
incidence variance in the opposite direction gives a total pencil
variance at most \(Cq^dBf_C\), with mean \(O(Bf_C)\).  Chebyshev's
inequality bounds its exceptional centers by \(Cn'e^{-2\xi}f_C\).
Sum this bound using \(\sum_Cf_C\le3\); the full-set and exceptional-set
indicator estimates have the same proof.  This argument pays for the
number of cells through their total mass, not their cardinality.

Poisson splitting makes the sample counts at different point--batch
pairs independent.  The average of \(\lambda_H\) over
\(H\in T\) is at most \(4\tau\).  Thus there is a fixed
\(T_0\subseteq T\) with \(|T_0|\ge|T|/2\) and \(\lambda_H\le8\tau\)
for every \(H\in T_0\).  For large \(q\), these hyperplanes lie in
\(\mathcal E\), and the product-size assumption gives
\(|T_0|\ge cqBe^{-b}\).
If \(Y\) counts the members of \(T_0\) empty in all batches, then
\[
 \EE Y\ge |T_0|e^{-8P\tau},\qquad 0\le Y\le |T_0|.
\]
The second bound, rather than independence of empty-hyperplane events,
yields
\begin{equation}
 \PP\{z\ge z_{\min}\}\ge\tfrac12e^{-8P\tau},
 \qquad z_{\min}=cqB e^{-b-8P\tau}.
 \label{eq:z-event}
\end{equation}
Here and below decrease the absolute constant \(c\) if necessary.

We may require at the same time
\begin{equation}
 \sum_{H\in Z}\lambda_H\le .001z.
 \label{eq:z-small-mass}
\end{equation}
The terms with \(\lambda_H\le10^{-4}\) contribute at most \(10^{-4}z\).
For the others, \eqref{eq:exceptional-mass} gives expected contribution
at most \(CqBe^{-10^{-4}P}\).  Markov's inequality at
\(.0009z_{\min}\) fails with probability at most
\(Ce^{b+8P\tau-10^{-4}P}=o(e^{-8P\tau})\), by
\eqref{eq:score-parameter-slack}.

For each resulting \(Z\), incidence variance shows that
\begin{equation}
 z_x\ge .8z/q
 \quad\hbox{outside at most }Cq^{d+1}/z\hbox{ ambient points}.
 \label{eq:z-ambient}
\end{equation}
On the training set,
\[
 \sum_{x\in S'}z_x=\frac{n'}q\sum_{H\in Z}\lambda_H\le .001n'z/q.
\]
Consequently at least \(.99n'\) training points satisfy \(z_x\le .1z/q\).


The exceptional part of the score equals \((1-b_0)^Rz_x\), with error in
absolute value at most \(b_0|\mathcal E\cap\{H:H\ni x\}|\).  For a member
of \(Z\) every factor is \(1-b_0\).  For any other exceptional hyperplane,
either an own hit makes its product zero or an outside-own hit supplies a
factor \(-b_0\); every other factor has magnitude at most one.  At a
regular pencil the error is at most \(Be^{-.91L}=o(z_{\min}/q)\), and
\((1-b_0)^R=1-o(1)\).  It remains to control the sum \(T_x\) over the
nonexceptional hyperplanes.  This is the only probabilistic issue left
in separating the two bounds on \(z_x\).

\subsection{The second moment on training points}
\label{subsec:score-second-moment}

Fix \(x\). The \emph{schedule} of a point or direction is the vector of
its sample counts in the $R$ batches. Condition first on \(x\) being
unsampled and then on the schedules of all points of
\(O_x\setminus\{x\}\).  The own-empty
multipliers in \eqref{eq:poisson-score} are now fixed zeros or ones.
All remaining point counts are still independent Poisson variables.
Group them by their radial lines \(\ell\) through \(x\), with strengths
\[
 a_\ell=\frac q{n'}|\ell\cap(S'\setminus(O_x\cup\{x\}))|.
\]
A line has a hit in a specified batch with probability \(1-e^{-La_\ell}\).
For a typical hyperplane put
\[
 \lambda'_H=\sum_{\ell\subset H}a_\ell,\qquad
 \delta'_H=\lambda'_H-(1-f),\qquad
 \theta_{HH'}=\sum_{\ell\subset H\cap H'}a_\ell.
\]
The union defining \(O_x\) uses two cells minus their intersection.
The exceptional thresholds therefore give
\begin{equation}
 \lambda'_H\ge .75,\qquad |\delta'_H|\le .17,\qquad
 D_x:=\sum_{H\text{ typical}}(\delta'_H)^2\le CBe^\xi
 \label{eq:typical-masses}
\end{equation}
at a regular pencil.  The last bound follows by summing the relevant
clipped squares.  Removing \(x\) contributes at most
\(O(q^{d-1}(q/n')^2)=O(B)\), since \(n'\ge n/4>25qP\).
Also \(\lambda'_H\le\lambda_H\le1.1\), so the geometric input applies.

Temporarily omit the fixed own multiplier and write the one-batch term as
\[
 X_H=\mathbf1\{H\text{ is outside-own empty}\}-b_0.
\]
Then
\[
 u_H:=\EE X_H=e^{-L\lambda'_H}-b_0,
 \qquad |u_H|\le Le^{-.75L}|\delta'_H|,
\]
and for distinct \(H,H'\),
\[
 \EE X_HX_{H'}=u_Hu_{H'}+
 e^{-L(\lambda'_H+\lambda'_{H'})}(e^{L\theta_{HH'}}-1).
\]
Since \(\lambda'_H+\lambda'_{H'}-\theta_{HH'}\ge .75\),
\begin{equation}
 |\EE X_HX_{H'}|\le L^2e^{-.75L}
       (|\delta'_H\delta'_{H'}|+\theta_{HH'}),
 \qquad \EE X_H^2\le2e^{-.75L}.
 \label{eq:one-batch-moments}
\end{equation}
The fixed multipliers can only reduce the absolute bounds.
More precisely, after the own schedules are fixed, the full product
for a hyperplane is either identically zero or has all its own-empty
multipliers equal to one.  Thus they do not alter the signed
cancellation used below.

In dimension two the overlaps vanish.  In higher dimensions,
\eqref{eq:score-geometric-inputs} and dyadic summation give
\[
 \sum_{H\ne H'}\theta_{HH'}^R
 \le B^2\exp(2\chi+O(R+\log\sigma)).
\]
For the deviation terms, \(R\ge2\) and \(|\delta'_H|<1\) give
\(\sum_{H,H'}|\delta'_H\delta'_{H'}|^R\le D_x^2\).
Finally \(Q_{d-1}=O(B^2)\), because
\(n'\le n\le Cq^{(d+1)/2}\).  Summing
\eqref{eq:one-batch-moments} to the \(R\)-th power and using
\eqref{eq:score-parameter-slack} proves
\begin{equation}
 \EE T_x^2\le B^2e^{-.6P}.
 \label{eq:validation-second-moment}
\end{equation}
This estimate is uniform in the conditioned own schedules.  It holds
at all centers except the \(o(n)\) pencil and geometric exceptions;
it uses none of the additional ambient exclusions.

In dimension two, different pencil lines have disjoint outside-\(x\)
samples.  Thus their score terms are independent under the same
conditioning.  Their variance sum is at most \(Cqe^{-.6P}\), and their
mean is at most \(Be^{-cP}\), by the one-batch mean bound and \(D_x\).
Bernstein's inequality gives
\begin{equation}
 \PP\{|T_x|>z_{\min}/(10q)\}\le q^{-50}.
 \label{eq:score-dimension-two}
\end{equation}
Indeed \(B\ge c\sqrt q\), \(P=o(\sigma)\), and \(e^{cP}\gg\sigma\).
This settles the ambient tail in dimension two.  The next subsection
treats dimensions three and four, where pencil tests share radial lines.

\subsection{An ambient moment bound}
\label{subsec:score-high-moment}

Choose the even moment order \(p\) above, and put \(K_c=200\), \(J=5000\).
If \(g_1>\chi_0\), let \(\mathcal G_x\) be the event that every
outside-own radial line hits in fewer than \(J\) batches.  Otherwise
\(\mathcal G_x\) is the whole sample space.  In the high case a line of
strength \(a\) violates the condition with probability at most
\((2Pa)^J\).  The global geometric line bound and the maximum strength
\(Ce^{-g_1}\) give, after summing the dyads above \(e^{-.05g_0}\),
\[
 \sum_x\PP(\mathcal G_x^c)
 \le Cq^2P^J e^{-(J-2)g_1+O(J)+O(\log\sigma)}+o(n)\le n.
\]
For the lower dyads, the elementary bound of \(Cq^7\) center--line pairs
in dimensions at most four contributes at most
\(Cq^7(CPe^{-.05g_0})^J=o(n)\); here \(g_0>.5\sigma\).
These are probabilities under the fixed-point conditioning, uniformly
in the own schedules.

We prove, at every remaining ambient center,
\begin{equation}
 \EE\bigl[T_x^p\mathbf1_{\mathcal G_x}\bigr]
 \le B^p e^{-.1pP}.
 \label{eq:ambient-high-moment}
\end{equation}
Expand the power into ordered \(p\)-tuples of typical hyperplanes.
Join two distinct tuple hyperplanes when their overlap exceeds
\(t=1/(100p)\).  A component is a connected component of this graph.
A simple singleton is an isolated hyperplane occurring in exactly one
tuple position. A positive-strength radial direction is \emph{external}
if it is contained in hyperplanes from two different components.
A batch \emph{touches} a hyperplane when it has a hit on an external
direction contained in that hyperplane. Thus sharing a direction always
means containing that radial line, not merely intersecting it at $x$.  The total external mass in any tuple hyperplane is at
most \((p-1)t<.01\): sum its overlap with tuple hyperplanes in other
components.

Condition on all external schedules.  Internal schedules belonging to
different components are independent.  Directions contained in no tuple
hyperplane contribute an independent truncation probability at most one,
which can be omitted from the absolute expectation bound.  The remaining
truncation condition splits into its external condition and one internal
condition per component.
The internal mass in every hyperplane is at least \(.74\).
In a batch with no external hit the expected absolute value of its factor is at
most \(2e^{-.74L}\).  If there is an external hit the factor has magnitude
\(b_0\le e^{-.92L}\).  A component containing two distinct hyperplanes,
or a singleton repeated in the tuple, consequently costs at most
\begin{equation}
 e^{-.7P}.
 \label{eq:nonsimple-attenuation}
\end{equation}
To see this, keep only the absolute factors of one of its hyperplanes;
all other absolute factors are at most one.  Drop its internal
truncation and multiply the one-batch bounds.  This bound is uniform in
the external schedules.

For a simple singleton \(H\), let \(\theta_{\max}\) be its largest overlap
with another tuple hyperplane, or zero if there is none.  The external
mass is at most \(p\theta_{\max}\le .01\).  In an untouched batch its
signed mean has magnitude at most
\(Le^{-.74L}(|\delta'_H|+p\theta_{\max})\); a touched batch has factor
\(-b_0\).  If fewer than half the batches are touched, the product of
means is therefore bounded in absolute value by
\begin{equation}
 e^{-.68P}\bigl(|\delta'_H|^{R/2}+
                         (p\theta_{\max})^{K_c}\bigr).
 \label{eq:simple-signed-attenuation}
\end{equation}
Here \((CL)^R=e^{o(P)}\), and \(R/2\ge K_c\).
If at least half are touched, the bound is instead
\(e^{-.68P}\mathbf1_{\mathrm{BAD}(H)}\), where
\(\mathrm{BAD}(H)\) denotes this condition on the external schedules.

Internal truncation adds to \eqref{eq:simple-signed-attenuation} at most
\begin{equation}
 e^{-.68P}\sum_{\ell\subset H}(Pa_\ell)^J.
 \label{eq:truncation-error}
\end{equation}
For a direct verification, union over a line with hits in \(J\) specified
batches.  Each forced batch has score factor \(-b_0\) and probability at
most \(La_\ell\).  Every remaining batch has absolute expectation at
most \(2e^{-.74L}\).  Summing the choices of forced batches gives
\((Pa_\ell)^J\) while retaining the displayed attenuation.  Thus no
independence conditional on the truncation event is being asserted.

We can now designate each simple singleton as a deviation, shift,
truncation-error or BAD term according to these bounds.  The first three
have respective weights
\[
 |\delta'_H|^{R/2},\qquad (p\theta_{\max})^{K_c},\qquad
 \sum_{\ell\subset H}(Pa_\ell)^J;
\]
the last has its BAD indicator.  The truncation-error designation is absent
when \(g_1\le\chi_0\), since \(\mathcal G_x\) is then the whole space.
Every simple position retains
\(e^{-.68P}\), including positions whose additional weight will be
dropped in the counting below.

\paragraph{Counting components and the first three designations.}
Normalize all counts by \(B\) for each tuple position.
Identifying repeated positions and choosing orders or component trees
costs \(p^{O(p)}\).  Begin each component with at least two distinct
hyperplanes by an ordered edge.  Its normalized count is at most
\(e^{2\chi}(C/t)^{K_c}\).  Every subsequent distinct vertex along a tree
costs at most \(\Delta_t/B\le e^{-3P}\); a repeated position costs
\(1/B\).  A repeated singleton has normalized cost \(O(1)\) for its
first two positions, since \(Q_{d-1}=O(B^2)\), and \(1/B\) thereafter.

Encode all deviation and truncation-error singletons next.  The
normalized sum of deviation weights is at most \(D_x/B\le Ce^\xi\).
For the error weights, the number of pencil hyperplanes containing a
fixed radial line is at most \(Cq^2\).  In the high case their normalized
sum is consequently bounded by
\[
 \frac{Cq^2}{B}\sum_\ell(Pa_\ell)^J\le e^{2\chi}.
\]
Indeed use the local line-count estimate dyadically, the identity
\((q^2/B)q^{2-2v}=O(e^{-g_1})\), and the maximum strength \(Ce^{-g_1}\).
The resulting bound is
\(\exp(\chi-(J-99)g_1+J\log P+O(J+\log\sigma))\), which suffices because
\(g_1>\chi_0\gg\log P\).

For a positive shift weight choose a neighbor attaining the largest
overlap, and its dyad \(a\le\theta<2a\).  If the neighbor is already
encoded, the normalized weighted choice costs at most
\[
 (\Delta_a/B)(2pa)^{K_c}\le e^{2\chi}(2p)^{K_c}.
\]
Otherwise encode the ordered pair using \(E_a/B^2\); the same
cancellation holds.  The neighbor is then a shift or BAD singleton,
because the other types have already been encoded.  Drop its extra
weight, which is at most one, but retain its attenuation.  Each operation
encodes one or two new vertices, so this procedure terminates even when
the chosen-neighbor relation contains cycles.

Only unencoded BAD singletons remain.  We next give the complete
certificate count for their schedule requirements.

\paragraph{Certificates for the BAD singletons.}
For each external direction \(\ell\) and batch \(r\), the hit indicator
has probability \(1-e^{-La_\ell}\le La_\ell\).  These indicators are
independent.  We use a dictionary of directions, maintaining the
following invariant: every dictionary direction belongs to an already
encoded hyperplane and has at least one specified positive hit event.
When a new direction is first specified, pay for its newly specified
direction--batch events.  Later references to this dictionary entry
impose no new hit requirement.  In particular they never pay for the
same random event twice.

This is a union bound over static descriptions.  For each fixed tuple
and description, the retained positive requirements involve distinct
direction--batch pairs and have probability at most the product of their
\(La_\ell\)'s.  Absence requirements may be dropped.  The description
chosen for an actual realization can depend on that realization; no
conditional independence along this choice is needed.

\paragraph{Stage 1: anchors.}
First consider \(g_1\le\chi_0\).  If an unencoded vertex has a touched
direction shared with an encoded hyperplane \(H_0\), use it as an anchor.
For a new direction, the weighted sum of its direction and batch choices
inside \(H_0\) is at most
\begin{equation}
 \sum_{r=1}^R\sum_{\ell\subset H_0}La_\ell\le 2P.
 \label{eq:anchor-weight}
\end{equation}
For an old direction use a dictionary reference and drop its hit
requirement.  The anchored hyperplane can be chosen in at most
\(Cq^{d-2}\) ways, giving normalized cost
\(Cq^{d-2}/B\le Cq^{v-1}\le Ce^{\chi_0}\).
Encode it and add any new direction.

If \(g_1>\chi_0\), apply this step only when there are two geometrically
distinct touched directions shared with encoded hyperplanes.  Choose
the two anchors in succession using \eqref{eq:anchor-weight} or dictionary
references.  They span a projective plane, so the normalized number of
possible hyperplanes is at most
\(Cq^{d-3}/B\le Cq^{v-2}\).  The two encoded neighbors need not be
distinct, but the two directions must be.  If this step is unavailable,
there is at most one touched direction to the encoded set.  On
\(\mathcal G_x\) it accounts for fewer than \(J\) batches.

\paragraph{Stage 2: new pairs.}
Recheck the anchor steps after every operation.  If no anchor step
applies, and two unencoded vertices share hits in at least \(K_c\)
distinct batches on directions contained in no encoded hyperplane, encode that
ordered pair.  Choose an overlap dyad \(a\le\theta<2a\) and exactly
\(K_c\) of those batches, with one witnessing direction in each.
All these directions are outside the dictionary.  Their specified
direction--batch events are distinct even if a direction repeats.
For the fixed pair the sum of their probability weights is at most
\((2Pa)^{K_c}\).  Combined with the pair count this costs at most
\begin{equation}
 (E_a/B^2)(2Pa)^{K_c}\le e^{2\chi}(2P)^{K_c}.
 \label{eq:pair-certificate-cost}
\end{equation}
Add both vertices to the encoded set and the new directions to the
dictionary, and resume the checks.  Restricting new-direction choices to entries outside the
dictionary preserves independence.  Enlarging their positive weighted
sum to all directions is only an upper bound on that sum; it does not
turn repeated events into independent ones.

\paragraph{Stage 3: residual roots.}
Suppose none of these operations remains possible.  At each residual
vertex, discard its fewer than \(J\) batches touching the encoded set
in the high case; in the low case there are none.  Choose a witnessing
direction and another residual vertex for each remaining touched batch.
No neighbor occurs in \(K_c\) such batches, since that would allow the
pair operation.  No direction occurs in \(K_c\) such batches either:
fix any residual neighbor containing that direction to obtain the same
contradiction.  Greedily selecting a witness and deleting those with
either its neighbor or its direction gives at least
\[
 h_0=\left\lfloor\frac{R/2-J}{2K_c}\right\rfloor
\]
distinct neighbors on distinct directions for every residual vertex.

There is a root set containing at most one tenth of the residual
vertices such that every nonroot has two of these witness neighbors
among the roots.  To prove this, select roots independently with
probability \(1/20\).  A fixed witness list has fewer than two selected
members with probability at most \((1+h_0)e^{-ch_0}\).  Its union bound
over at most \(p\) vertices tends to zero, since \(R\gg\log p\).
Markov's inequality gives probability at most \(1/2\) that more than one
tenth are selected.  Thus both desired properties hold for some choice.
A nonempty residual set has at least \(h_0+1\) vertices, so there is no
small-set exception to this argument.

Encode the roots freely, dropping their BAD indicators.  Encode every
nonroot from two distinct directions to roots, as in the high-case
anchor step, using dictionary references whenever possible.  If there
are \(h\) residual vertices and \(r\le h/10\) roots, the aggregate
normalized geometric cost is at most
\begin{equation}
 C^h q^{rv+(h-r)(v-2)}
 =C^h q^{h(v-2)+2r}\le C'^h q^{-.3h},
 \label{eq:root-cancellation}
\end{equation}
because \(v\le1.5+O(1/\sigma)\).  The geometric restrictions remain valid
when a direction is already in the dictionary: reusing its name drops
only a hit requirement, not the requirement that it lie in the new
hyperplane.  In particular, two reused but distinct directions still
span a plane.  Their previously paid hit events are not charged again,
while the new hyperplane must still contain that plane.

At most \(CK_cp\) dictionary entries and positive hit requirements have
been recorded.  Orders, statuses, vertex names, batches, dictionary
references and dyads, together with the powers of \(P\) in
\eqref{eq:anchor-weight} and \eqref{eq:pair-certificate-cost}, cost at most
\begin{equation}
 (CpR\sigma P)^{CK_cp}=\exp(o(pP)).
 \label{eq:certificate-overhead}
\end{equation}
The non-overhead factors are at most \(e^{2\chi+O(1)}\) per operation;
\eqref{eq:root-cancellation} is favorable.  This completes the count of
all residual BAD requirements.

Some counting bounds enlarge the sum over valid descriptions by dropping
geometric constraints.  The extra configurations need not have the
original strong-component pattern: they are only positive summands in a
numerical upper bound.  Attenuation is applied to valid tuples before
their description sums are enlarged, not to these extra configurations.

\paragraph{Completing the moment estimate.}
All component attenuations were uniform in the external schedules.
They therefore remain valid for every static certificate.  Only now
integrate the external variables, dropping the external truncation
indicator before using their product law.  There is no assertion that
these variables are independent conditional on \(\mathcal G_x\).

A nonsimple component pays for its first two positions by
\(e^{-.7P+2\chi}\), apart from the overhead.  Each extra distinct position
has factor \(e^{-3P}\), and each repeat has factor \(1/B\le e^{-P}\) for
large \(q\).  Every simple position keeps \(e^{-.68P}\), with cost at most
\(e^{2\chi+o(P)}\) per position.  Deviation costs \(e^\xi\) are smaller.
After the overhead \eqref{eq:certificate-overhead}, these bounds give
\eqref{eq:ambient-high-moment}, with room to spare.  The moment order is
even, so Markov's inequality applies despite signed individual score
terms.  Since \(z_{\min}/(10q)=cB e^{-b-8P\tau}\), it gives
\begin{equation}
 \PP\{|T_x|>z_{\min}/(10q),\ \mathcal G_x\}
 \le \exp[-.1pP+p(b+8P\tau+O(1))]\le q^{-50}.
 \label{eq:ambient-score-tail}
\end{equation}

\subsection{Simultaneous capture and size}
\label{subsec:score-assembly}

We finish the proof of Lemma~\ref{lem:score-procedure} under the original
sampling law.  All fixed-point estimates above were conditional on the
point being unsampled and uniform in its own schedules.  Multiplying by
the probability that it is unsampled and then averaging the own
schedules makes them unconditional upper bounds for failures at
unsampled points.  Sampled points are included automatically.

The deterministic pencil and local geometry exceptions contain
\(o(n)\) points.  On the training set, the second moment
\eqref{eq:validation-second-moment} gives expected typical-score failures
at most \(n'e^{-cP}\).  Markov's inequality shows that more than \(.01n'\)
such failures has probability \(e^{-c'P}\), which is negligible relative
to \(e^{-8P\tau}\).  This validation estimate does not discard the
additional \(ne^{O(P)}\) ambient centers.

For the ambient size estimate, those additional centers may all be
included.  Beyond all static exceptions, \eqref{eq:ambient-score-tail},
\eqref{eq:score-dimension-two}, and the bound on
\(\sum_x\PP(\mathcal G_x^c)\) show that the expected number of dynamic
failures is \(O(n+1)\).  Markov's inequality permits at most \(ne^{CP}\)
of them, except with probability \(e^{-cP}\), again negligible relative
to \eqref{eq:z-event}.  The total Poisson sample count has mean \(qP\)
and exceeds \(2qP\) with probability \(e^{-cqP}\).

Intersect these events with \eqref{eq:z-event} and
\eqref{eq:z-small-mass}.  Their probability is at least
\(ce^{-8P\tau}\).  At least \(.99n'\) training points have
\(z_x\le .1z/q\); after the \(o(n)\) static exceptions and the validation
failures, a fixed positive fraction have score below \(.5z/q\) or are
sampled.  They belong to \(W\).  Outside the ambient exceptions,
\(z_x\ge .8z/q\), the exceptional-score error is \(o(z/q)\), and the
typical score has magnitude at most \(z/(10q)\).  Such unsampled points
have score above \(.5z/q\) and are excluded.

Finally, the \(Z\)-dependent ambient exception in
\eqref{eq:z-ambient} has size
\[
 Cq^{d+1}/z\le Cn'e^{b+8P\tau}\le ne^{CP}.
\]
The at most \(2qP\) sampled points, all other ambient exceptions and
dynamic failures satisfy the same size allowance.  Thus
\(|W|\le ne^{CP}\) and \(|W\cap S'|\ge cn'\) simultaneously on a set of
rows of probability at least \(\exp[-C(P\tau+1)]\), proving the lemma.
